\documentclass[10pt,twocolumn]{article}

\usepackage[T1]{fontenc}
\usepackage[utf8]{inputenc}
\usepackage[margin=0.7in]{geometry}
\setlength{\columnsep}{0.24in}
\usepackage{amsmath,amssymb}
\usepackage{graphicx}
\usepackage{enumitem}
\usepackage[colorlinks=true,linkcolor=blue,citecolor=blue,urlcolor=blue]{hyperref}
\usepackage{titlesec}
\usepackage{stfloats}
\usepackage{balance}
\setcounter{topnumber}{2}
\setcounter{bottomnumber}{2}
\setcounter{totalnumber}{4}
\renewcommand{\topfraction}{0.92}
\renewcommand{\bottomfraction}{0.75}
\renewcommand{\textfraction}{0.06}
\renewcommand{\floatpagefraction}{0.80}

\titlespacing*{\section}{0pt}{0.9ex}{0.4ex}
\titlespacing*{\subsection}{0pt}{0.7ex}{0.3ex}
\titleformat*{\section}{\large\bfseries}
\titleformat*{\subsection}{\normalsize\bfseries}
\setlist{nosep,leftmargin=1.2em}
\setlength{\parskip}{0.2ex}
\setlength{\abovecaptionskip}{4pt}
\setlength{\belowcaptionskip}{0pt}
\raggedbottom
\setlength{\floatsep}{8pt plus 1pt minus 1pt}
\setlength{\textfloatsep}{10pt plus 1pt minus 1pt}
\setlength{\intextsep}{8pt plus 1pt minus 1pt}
\setlength{\dblfloatsep}{8pt plus 1pt minus 1pt}
\setlength{\dbltextfloatsep}{10pt plus 1pt minus 1pt}

\title{\vspace{-2.4em}\large A Continuous Attractor in Mouse Motor Cortex,
and a Causal Readout for Closed-Loop Intervention}
\author{
Samridh Tiwari$^{1,2}$ \quad Om Kolhe$^{1}$ \quad Krishna Jayant$^{1}$\\[2pt]
\small $^{1}$Nano Neurotechnology Laboratory, Weldon School of Biomedical
Engineering, Purdue University, West Lafayette, IN 47907\\
\small $^{2}$Indian Institute of Technology Bombay, Mumbai 400076, India\\[1pt]
\small \texttt{tiwar133@purdue.edu} $\cdot$ \texttt{samridhtiwari2@gmail.com}}
\date{\vspace{-1.6em}}

\begin{document}
\makeatletter
\twocolumn[
  \begin{@twocolumnfalse}
  \maketitle
  \vspace{-2.4em}
  \begin{abstract}\noindent\small
A motor plan must be both \emph{held} against noise and \emph{changed} when new
information arrives. We fit LFADS to spiking activity from two mice (9 sessions,
39--118 neurons) in a cued lever-pull task. It beat GPFA at matched latent
dimensionality on entirely held-out neurons in 8 of 8 fitted sessions (median
$4.1\times$), and a lagged linear readout of its state recovered the continuous
lever trajectory on held-out trials ($R^2=0.47$--$0.77$). Run autonomously from
one inferred initial condition with all input removed, the learned rule
reproduced 94\% of pre-cue single-trial latent variance against 2\% for a
condition-mean predictor. Linearising the generator at 512 fixed points revealed
a \textbf{continuous attractor}: $\sim$33 marginally stable directions, $\sim$95
relaxing within $\sim$84~ms, and essentially none unstable, replicated across
five models per animal and both animals. Perturbations along the manifold
retained 75\% of their magnitude after 1.2~s against 8\% off-manifold. Because
every such estimate is acausal, we built a \textbf{causal real-time detector} of
impending movement from unsorted multiunit activity: streaming live it flagged
movement a median of \textbf{80~ms before lever motion} at \textbf{under 1~ms
compute per 20~ms bin}, beating a circular-shift null by 6.7~s.d.\ in precision.
Zero-shot cross-day transfer collapsed to chance (AUC 0.485) but was restored by
latent alignment (0.740 against a 0.768 within-day ceiling), localising the
failure to a rotation of readout coordinates rather than a change in dynamics.
  \end{abstract}
  \vspace{1.0em}
  \end{@twocolumnfalse}
]
\makeatother

\section{Introduction and Background}

A voluntary movement is preceded by a preparatory period in which motor cortical
activity evolves without producing muscle output \cite{churchland2012}.
This activity occupies a low-dimensional subspace with reproducible trajectories, and trajectories rather than single-neuron tuning are the
right level of description \cite{churchland2012}.

This creates a tension: the plan must be \emph{stable}, or noise corrupts it,
yet \emph{changeable}, or a cue cannot move it. A classical resolution is a
\textbf{continuous attractor}, a connected manifold of marginally stable
states \cite{seung1996,zhang1996}. Displacements along it neither grow nor decay,
so information is held; displacements off it are corrected, so noise is
rejected. Continuous attractors are established for head-direction and
oculomotor integration \cite{zhang1996}, but evidence in motor
preparation is sparser, because the analysis needs a \emph{generative model of
the dynamics}, not a description of trajectories. LFADS \cite{pandarinath2018}
supplies one: a sequential variational autoencoder whose decoder is trained to
generate the observed spikes. Once fit, the generator is a dynamical system that
can be interrogated directly: fixed points located numerically, Jacobians
computed exactly, stability read from the eigenspectrum \cite{sussillo2013}.

Our motivation is interventional: to nudge a plan toward a successful
configuration, a closed-loop system needs a target state, a model of what a
perturbation does, and a real-time estimate of where the circuit is. Part~A
addresses the first two; Part~B the third.

\section{Methods}

\subsection{Task, recordings, preprocessing}

Two mice (\texttt{61199SOM}, \texttt{61198SOM}) performed a cued forelimb
lever-pull task; a tone signalled trial onset and a successful pull (``Hit'')
was rewarded with water. Spiking was recorded from forelimb motor cortex across
\textbf{9 sessions} (4 and 5 per animal), \textbf{39--118} isolated units per
session, trials aligned to cue onset from $-1.5$ to $+1.5$~s. LFADS was fit on
the 8 sessions in the stitched archive: \texttt{61199SOM}, 1163 trials (844
Hit); \texttt{61198SOM}, 956 trials (595 Hit). Counts were rebinned to
\textbf{20~ms} (150 bins/trial), split chronologically 80/20; no smoothing
preceded fitting. Part~B sessions used a 64-channel silicon probe at 20~kHz
(Intan RHD2000) with analog lever position and TTL cue/reward lines on one
clock; two dual-probe sessions recorded M1 and M2 simultaneously. \textbf{No
spike sorting was performed in Part~B}. The detector uses threshold-crossing
multiunit activity, preserving population structure while removing a step that
cannot be done reliably online \cite{trautmann2019}.

\subsection{Models, selection, evaluation}

\textbf{LFADS} \cite{pandarinath2018} was fit with \texttt{lfads-torch}
\cite{sedler2023}, multi-session with per-session read-in/read-out and a shared
generator: 128-unit generator, 15 factors, controller output 4,
initial-condition dimension 64, dropout 0.1, batch 32, lr $4\times10^{-3}$,
Poisson observations. \textbf{GPFA} \cite{yu2009} was fit at matched
dimensionality (\texttt{xDim}=15). A sweep over KL and L2 scales gave 10
configurations per animal, ranked by validation reconstruction loss; we then
computed an \textbf{autonomy score} (correlation between the generator run
forward from its inferred initial condition with input zeroed, and the full
model) and selected \texttt{cfg06}, not the reconstruction winner
(\S\ref{sec:sel}).

\textbf{Held-out-neuron co-smoothing} \cite{pei2021}: neurons and trials each
split into 5 folds; factors inferred from retained neurons, a Poisson GLM fit to
each held-out neuron, log-likelihood scored on held-out
neurons~$\times$~held-out trials as bits/spike against that neuron's mean rate.
Both models used an identical target and GLM, isolating latent quality.
\textbf{Kinematic decoding} used a lagged ridge readout ($-60$ to $+160$~ms of
generator state), 5-fold cross-validated across trials, against the
training-set average pull as baseline.

\textbf{Fixed points} were located by minimising
$\|f(\mathbf{g},\mathbf{0})-\mathbf{g}\|^2$ with Adam (lr $2\times10^{-2}$,
cosine annealing to $10^{-5}$, 12{,}000 iterations) from 512 initialisations
drawn from observed states with noise ($\sigma=0.05$) \cite{sussillo2013}, kept if relative displacement fell below
$10^{-3}$. Jacobians came from exact autodifferentiation of the trained GRU
cell; directions were classed \textbf{marginal} ($\big||\lambda|-1\big|<0.01$),
\textbf{stable} ($|\lambda|<0.99$) or \textbf{unstable} ($|\lambda|>1.01$),
$\tau=-\Delta t/\ln|\lambda|$. For \textbf{perturbation analysis} the full
nonlinear generator was displaced from a fixed point along a marginal or random
non-marginal eigendirection and released; the eigendecomposition selects the
direction only, all simulation uses the trained network.

\subsection{Causal real-time detection}

Every estimate above depends on $g_0$, which LFADS infers with a bidirectional
encoder spanning the whole trial (tip~3); none can be produced online. Part~B
asks a narrower question: can the \textbf{trigger} a closed-loop system needs be
obtained causally, within the latency the behaviour allows?

A discriminative detector runs single-pass over the raw stream: common-median
referencing across 64 channels; a causal 3rd-order Butterworth high-pass at
300~Hz; a per-channel adaptive threshold at $4\times$ a running
median-absolute-deviation estimate; crossings with 1~ms refractory; accumulation
into \textbf{20~ms bins}; z-scoring; and concatenation of the \textbf{12 most
recent bins (240~ms)} into a $64\times12=768$-D vector. Filter and baseline
states carry across file boundaries and never reset, so output on a recording is
identical to the live stream. An alarm fires when the projected score exceeds
$\theta$ for two consecutive bins (40~ms dwell). Two readouts were compared: a
regularised Fisher discriminant, and a PCA-whitened cosine similarity to a
movement template, bounded to $[-1,1]$.

Training uses lever-derived labels only; cue timing is never supplied, so the
detector is agnostic to whether a movement is cued or self-initiated. Positives
are bins in $[-200,0]$~ms before a pull reaching the hardware reward threshold
(tip~4); negatives are bins $>2$~s from any pull at rest. Cross-validation is
\textbf{time-blocked}. Baselines: an unsupervised CUSUM \cite{page1954} ramp
detector on the first principal component of the same features; and a
\textbf{circular-shift null} rotating alarm times on the session clock,
preserving alarm rate and periodicity while destroying only alignment to
movement. This matters because the animal pulls every $\sim$3.7~s. Cross-day
transfer was tested zero-shot, then after Procrustes registration in the 15-D
latent space using pull-triggered averages as anchors
\cite{gallego2020,degenhart2020}.

\subsection{Sample, Code, and Data Storage}

Analysis code, trained model checkpoints, the derived numbers behind every
reported value, and the editable source for every figure are deposited at
\url{https://github.com/Samridh-Tiwari/surf2026-motor-dynamics}. Each figure is
produced by a committed script that also emits a \texttt{*\_stats.json} of the
values it reports, and the \texttt{README} maps every reported number to the
script that generates it. The \textbf{pre-processed} data (raw Intan
\texttt{.rhd} recordings with their behavioural channels, and the stitched spike
archive) and the \textbf{post-processed} data (rebinned model inputs, LFADS
per-session outputs, binned dual-probe arrays) exceed the repository size limit
and are provided on request to the corresponding author. The committed checkpoints are sufficient to regenerate
every Part~A result without retraining.

\textbf{Labelling scheme.} Animals \texttt{61199SOM} and \texttt{61198SOM};
sessions \texttt{Day2}--\texttt{Day7} within an animal; the stitched archive
indexes these as \texttt{dataset\_01}--\texttt{08}; LFADS configurations
\texttt{cfg01}--\texttt{cfg10}, with \texttt{cfg06} selected throughout.
Environment: conda \texttt{lfads-torch} with PyTorch/CUDA (Part~A) and MATLAB
R2026a (Part~B); all stochastic procedures are seeded, so the reported numbers
reproduce exactly.

\subsection{Reference Methods}

LFADS: \url{https://github.com/arsedler9/lfads-torch} \cite{sedler2023}. GPFA:
\url{https://github.com/NeuralEnsemble/elephant} \cite{yu2009}. Fixed-point
location follows \cite{sussillo2013}, co-smoothing the Neural Latents Benchmark
protocol \cite{pei2021}, and sorting-free population analysis
\cite{trautmann2019}.

\subsection{Insider Tips}

\begin{enumerate}
\item The LFADS factor read-out is \texttt{KernelNormalizedLinear}, which
L2-normalises each weight row; a plain matmul with the stored weights does
\textbf{not} reproduce the saved factors. Apply
\texttt{F.normalize(W,p=2,dim=1)} first.
\item \texttt{gen\_init} is the 128-D generator state \emph{after} the
\texttt{ic\_to\_g0} projection, not the 64-D initial-condition sample; saved
\texttt{gen\_states} are posterior means over 50 samples, so re-simulating from
the mean does not reproduce them exactly.
\item The initial-condition encoder is \textbf{bidirectional over the whole
trial}, so anything computed from $g_0$ or the factors is acausal, which motivates Part~B.
\item \textbf{The lever threshold triggering reward is fixed hardware and must be
recovered as an absolute voltage, not a session statistic} (2.8197 and 2.8188~V
on two days). Thresholds from the median absolute deviation or amplitude
percentiles drift and silently admit sub-threshold micro-movements; on one
session this inflated the pull count from 864 to 2702. Recover it from the lower
edge of the peak lever voltage preceding each reward.
\item Behavioural events live in the raw \texttt{.rhd} files
(\texttt{ANALOG-IN-2}=lever, \texttt{DIGITAL-IN-04}=cue,
\texttt{DIGITAL-IN-03}=reward), not the exported workspace; the default Intan
reader discards analog board channels. Carry causal filter and baseline state
across file boundaries, since re-initialising per file produces a transient the
detector reads as movement.
\item Kinematic decoding needs a \textbf{lagged} readout: an instantaneous one
scores $R^2=0.32$ where a lagged one reaches 0.77, because motor cortex leads
movement.
\end{enumerate}

\subsection{Deviations from Standard Methodology}

\begin{enumerate}
\item \textbf{Model selection} on autonomy subject to near-best reconstruction,
because reconstruction is insensitive to the dynamics/input trade-off
(\S\ref{sec:sel}).
\item \textbf{Perturbation actuator.} Time-varying LQR through the model's 4
inferred input channels yields 0.1\% error reduction, because those channels have
negligible authority over a 128-D state. We therefore perturb the state
directly, reflecting a physical intervention.
\item \textbf{GPFA dimensionality} fixed at 15 to match LFADS rather than
cross-validated (stored \texttt{cvLL} is empty): a matched comparison, not one
at each method's own optimum.
\item \textbf{Readout selection} favours the bounded cosine over the Fisher
discriminant despite indistinguishable held-out accuracy, on robustness grounds.
\item \textbf{Cross-validation} is time-blocked, not random $k$-fold; random
folds leak between adjacent trials and inflated one cross-area difference by
$\sim$40\%.
\end{enumerate}

\subsection{Methods Planned for Completion}

In progress: \textbf{causal state estimation} (a forward-only encoder over a
fixed past window, replacing the bidirectional one, to obtain a manifold
coordinate online rather than a binary trigger), and an \textbf{optogenetic
protocol} stimulating along versus across the manifold, gated by the validated
trigger.

\section{Results}

\subsection{Latent quality}

PCA of trial-averaged activity showed $\sim$3 components capturing 95\% of
variance with no increase as neuron count rose (Spearman $r=-0.16$, $p=0.70$).
Since smoothing inflates this, we use it only to justify a 15-D latent space
chosen with headroom rather than fitted.

LFADS achieved higher held-out-neuron predictive information in \textbf{8 of 8
fitted sessions} (0.303 vs 0.080 bits/spike; median ratio $\mathbf{4.1\times}$,
range 2.6--5.9$\times$). A two-sided Wilcoxon signed-rank test gives $p=0.008$,
but that is the \emph{smallest value attainable with 8 paired samples}
($2/2^8=0.0078$), so it confirms only that every session favoured LFADS and says
nothing about effect size; the effect size is the ratio. Held-out neurons are
never seen by either model, so this cannot be satisfied by memorisation. We hold
out \emph{neurons} rather than \emph{trials} deliberately: LFADS infers
per-trial initial conditions and inputs, so it fits trials nearly by
construction.

A lagged linear readout of the generator state also recovers the
\emph{continuous lever trajectory} on held-out trials (Fig.~\ref{fig:lever}), $R^2=0.47$--$0.77$ across
four sessions, against 0.26--0.68 for predicting each held-out trial with the
training-set average pull, a strong baseline given how
stereotyped the movement is. The model was fit to spikes alone and never saw the lever.

\begin{figure}[!htb]\centering
\includegraphics[width=\columnwidth]{lever_decode_traces.png}
\caption{A model fit to spikes alone predicts the lever. Five held-out trials:
real lever position (grey) and the lagged linear readout of generator state
(green), each trial excluded from the fold that predicts it. Across four
sessions $R^2=0.47$--$0.77$, against 0.26--0.68 for the training-set average
pull.}
\label{fig:lever}
\end{figure}

\subsection{Reconstruction loss selects against interpretable dynamics}
\label{sec:sel}

The sweep's reconstruction winner (\texttt{cfg01}) proved unsuitable for
dynamical analysis: run autonomously its generator reproduced only 48\% of the
state trajectory (45\% in mouse~2), the rest being carried by inferred inputs.
\texttt{cfg06}, at 2.9\% worse reconstruction, was \textbf{97\% autonomous}
(86\% in mouse~2), and the ordering held in both animals. The two differ in
exactly one respect: \texttt{kl\_co\_scale} is $10\times$ smaller in
\texttt{cfg01}. A weaker penalty on controller outputs lets the model explain
data with injected input, which reconstruction loss rewards equally to
explaining it with dynamics. \textbf{Reconstruction loss is therefore not a
valid selection criterion when the claim concerns intrinsic dynamics.}

Iterating \texttt{cfg06} from one inferred initial condition with \textbf{all
input removed} reproduced 94\% (pre-cue) and 82\% (post-cue) of single-trial
latent variance. Three controls: a condition-mean predictor achieved
$R^2=0.02/0.03$, so trials are not stereotyped; held-out trials scored 0.96/0.78
against 0.93/0.82 on training trials, so there is no overfitting; a shuffled
initial condition gave $R^2<0$. Because $g_0$ is inferred from the whole trial
this concerns \textbf{compressibility}, not prediction.

\subsection{A continuous attractor}\label{sec:attractor}

All 512 seeds converged to distinct fixed points, itself the signature of a
connected manifold rather than isolated attractors. Across 60 sampled fixed
points, eigendecomposed Jacobians gave \textbf{33.2 marginal}, \textbf{94.7
stable} and \textbf{0.08 unstable} directions per fixed point, stable directions
relaxing at median $\tau=84$~ms (Fig.~\ref{fig:attr}). The absence of unstable
directions held in every model examined, and the structure replicated across
\textbf{five independently trained models per animal} (median marginal count 34,
range 20--46 in mouse~1; 36, range 8--50 in mouse~2). Marginal subspace
dimensionality varies with regularisation and sampling; its existence, and the
absence of unstable directions, do not. Mean successful-trial paths from
different sessions, and from the two animals' separately trained models, match
after rigid alignment in the 15-D factor space (shape difference 0.02--0.21
across sessions, 0.11 across animals; 1.00 is chance;
Fig.~\ref{fig:gen}).

\begin{figure}[!htb]\centering
\includegraphics[width=\columnwidth]{crossanimal.png}
\caption{Mean successful-trial path from the two animals, fit by separate models
with no shared parameters, after shift, scale and rotation only. Dots mark cue
onset. Compared in the model's full 15-dimensional output space the shape
difference is 0.11, against 0.99$\pm$0.01 for alignment to a time-scrambled
target.}
\label{fig:gen}
\end{figure}

\begin{figure*}[!tb]\centering
\includegraphics[width=0.88\textwidth]{fig3_attractor_spectrum.png}
\caption{The trained generator implements a continuous attractor.
\emph{Left}, exact Jacobian eigenvalues at converged fixed points against the
unit circle; the spectrum crowds the circle from inside with essentially nothing
outside. \emph{Centre}, sorted eigenvalue magnitude per fixed point (mean $\pm$
s.d.): a plateau of marginal modes at $|\lambda|\approx1$ gives way to a
decaying stable spectrum. \emph{Right}, the fixed-point manifold is
low-dimensional, comparable to the data. Selected model \texttt{cfg06},
\texttt{61199SOM}, 512 converged fixed points.}
\label{fig:attr}
\end{figure*}

\subsection{The cue as inferred input; the cost of intervening}

The model received spike trains only, never cue timing and never outcome, yet
the inferred input $\|\mathbf{u}_t\|$ was flat before the cue and rose sharply
after, peaking at \textbf{1.73$\times$ baseline at $+100$~ms} (mouse~1) and
\textbf{2.18$\times$ at $+140$~ms} (mouse~2), and larger on successful trials
($P(\text{Hit}>\text{Miss})=0.619$ and $0.813$; 2000-fold permutation null 95th
percentile 0.532). The modulation differs in size between animals, the
cue-locked transient does not.

Displacing the full nonlinear network along a marginal direction left
\textbf{75\%} of the displacement intact after 1.2~s; along a random
non-marginal direction, \textbf{8\%} survived. A transient intervention
therefore persists only if aimed within the manifold. The behaviourally relevant
axis lay preferentially within it (0.352 of its energy versus 0.254 expected;
$1.38\times$, exceeding all 400 random directions, $p<0.0025$). The effect is modest, since most of that direction remains off-manifold, so steering behaviour requires sustained rather than impulsive input. The geometry is nonetheless not arbitrary with respect to behaviour.

\subsection{A causal detector flags movement before it occurs}

On held-out data the detector reached \textbf{AUC 0.877}, catching 60\% of
movements at $\sim$5 false alarms per minute, with alarms a median of
\textbf{63~ms} early and \textbf{92\%} arriving before the movement ($n=199$).
The unsupervised CUSUM baseline on the same features did worse, so the
supervised projection, not the front end, carries the discrimination.

Deployed live, calibrating from the lever as the session ran, it flagged
movement a median of \textbf{80~ms} before lever motion, \textbf{90\%} of alarms
early and \textbf{71\%} followed by a real movement, of which \textbf{46
preceded self-initiated} and \textbf{20 cued} pulls, from a single template never
told which was which.

\textbf{Latency.} Median compute was \textbf{under 1~ms per 20~ms bin} on one
CPU core, under 5\% of budget, so the irreducible latency is structural: the
20~ms bin plus 40~ms dwell give a 60~ms floor. Against the time-resolved
analysis below, which finds movement decodable from $\sim$140~ms before onset,
$140-60=80$~ms is the predicted lead, and that is what the live system delivers.

\textbf{Readout choice.} Both readouts achieved indistinguishable held-out AUC
(0.83--0.85) but differ in the tails: the Fisher discriminant is unbounded and
produced rest-period excursions of 6--23~s.d., intermittently straddling its
threshold live and firing on artefacts. \textbf{At equal nominal accuracy,
bounded scores should be preferred for closed-loop use}, a distinction invisible
to offline AUC.

\textbf{Base-rate correction.} The animal moves roughly every 3.7~s, so raw hit
rate is uninterpretable. Against a circular-shift null (5000 permutations) the
detector achieved precision \textbf{74.2\% against 39.3\%$\pm$5.2 expected}
($z=+6.7$) and recall \textbf{35.3\% against 16.8\%$\pm$2.3} ($z=+8.2$), beating
99.9\% and 100\% of permutations, with lead times $6.3\times$ more concentrated
in the $[0,+150]$~ms window than uniform.

\subsection{Cross-day transfer; what a second area contributes}

Applied zero-shot to the next day the model collapsed to \textbf{AUC 0.485},
i.e. chance, with false alarms rising to 21.8/min. After Procrustes alignment in
the 15-D latent space it reached \textbf{AUC 0.740} against a within-day ceiling
of 0.768. The information survives; what changes is the coordinate system in
which electrodes report it, consistent with stable latent dynamics beneath
unstable single-channel readouts \cite{gallego2020,degenhart2020}. Seeding
calibration with prior-day data nearly doubled early-session recall
(32\%$\to$56\%). Time-resolved decoding of movement against matched rest became
reliable \textbf{140--260~ms before onset} at peak AUC 0.85, which sets the
ceiling on any causal system's lead time.

The rig already carries two probes, so the question asked of the premotor
recording was not whether it buys lead time but whether M1 and M2 together
\emph{distinguish self-initiated from cue-initiated movements}. On lead time it
contributes nothing: the two population signals cross-correlate at \textbf{0~ms
lag} ($r=0.98$--0.99, both days) with indistinguishable ramp latencies, and
combining probes detected no earlier than M1 alone ($-120$ vs $-140$~ms), consistent with the delay-line organisation of inter-areal motor sequences
reported in this laboratory \cite{khan_delaylines}. M2's apparently higher
balanced accuracy (75.7\% vs 71.9\%) failed every control: it reversed for long
cue-to-movement delays ($+2.6\to-4.0$), vanished under reward-validated labels
(81.8 vs 81.5), and was absent for self-initiated movements ($n=447$, n.s.).
That localises M2's extra signal to the \textbf{cue-evoked response}, which is
the component separating cued from self-initiated trials and therefore the
usable answer to the question posed, rather than an earlier intention signal.

\section{Discussion, Conclusions, and Future Directions}

The continuous attractor resolves the tension of \S1 by geometry rather than tuning, holding information along the manifold and rejecting noise across it. This connects motor preparation to a mechanism established for head-direction and
oculomotor integration \cite{seung1996,zhang1996}.

\textbf{For intervention} this is a design constraint, not a prescription. A
single perturbation persists only if aimed within the manifold (75\% vs 8\%),
and the behavioural axis is tilted toward it but far from contained by it, so a
practical controller needs sustained input with a substantial off-manifold
component. The negative LQR result is informative: the model's inferred input
channels are too low-dimensional to steer a 128-D state, so a physical
intervention must access the state more directly than the model's
parameterisation suggests, motivating the nanoscale and optical interfaces
developed in this laboratory \cite{jayant2019,gonzales2025}.

\textbf{What Part~B does and does not close.} The \emph{timing} requirement is
met, and the operative engineering constraint identified: the readout rotates
between days, so a deployed system needs a cheap per-session realignment. It
does \textbf{not} close the state-estimation gap. The detector is discriminative,
answering ``is a movement imminent?'', not ``where on the manifold is the circuit
now?'', and aiming a perturbation needs the latter.

\textbf{Limitations.} Part~A state estimates are \textbf{acausal} and cannot run
online, and the detector is discriminative rather than generative. Marginal
subspace dimensionality is regularisation- and sampling-dependent (8--50), so the
\emph{count} should not be over-interpreted though the \emph{existence} is
robust. Live validation rests on one session in one animal and the cross-area
comparison on two.

\textbf{Future directions.} Causal or streaming LFADS-class inference, to obtain
a manifold coordinate online rather than a binary trigger; optogenetic
perturbation along versus across the manifold, which the steering result makes
falsifiable and the validated trigger can gate; and whether manifold
dimensionality changes over learning, testable in the existing multi-day
recordings.

\section*{Acknowledgements}
\footnotesize
I thank \textbf{Om Kolhe} for mentorship throughout this project and for the
dual-probe recordings analysed here, \textbf{Hammad Khan} and \textbf{Shivam} for
technical guidance, and \textbf{Prof.\ Krishna Jayant} for hosting this work and
directing its scope. All recordings were performed by members of the laboratory;
my contribution was computational. Thanks also to \textbf{Monish Lokhande}.
Supported by the Purdue SURF program.

\balance
\begingroup\tiny
\begin{thebibliography}{99}\itemsep0pt\parsep0pt\parskip0pt
\bibitem{churchland2012} Churchland, M. M., et al. \emph{Nature} \textbf{487},
51 (2012).
\bibitem{seung1996} Seung, H. S. \emph{PNAS} \textbf{93}, 13339 (1996).
\bibitem{zhang1996} Zhang, K. \emph{J. Neurosci.} \textbf{16}, 2112 (1996).
\bibitem{pandarinath2018} Pandarinath, C., et al. \emph{Nat. Methods}
\textbf{15}, 805 (2018).
\bibitem{sussillo2013} Sussillo, D. \& Barak, O. \emph{Neural Comput.}
\textbf{25}, 626 (2013).
\bibitem{sedler2023} Sedler, A. R. \& Pandarinath, C. \emph{arXiv:}2309.01230
(2023).
\bibitem{yu2009} Yu, B. M., et al. \emph{J. Neurophysiol.} \textbf{102}, 614
(2009).
\bibitem{pei2021} Pei, F., et al. \emph{arXiv:}2109.04463 (2021).
\bibitem{trautmann2019} Trautmann, E. M., et al. \emph{Neuron} \textbf{103}, 292
(2019).
\bibitem{gallego2020} Gallego, J. A., et al. \emph{Nat. Neurosci.} \textbf{23},
260 (2020).
\bibitem{degenhart2020} Degenhart, A. D., et al. \emph{Nat. Biomed. Eng.}
\textbf{4}, 672 (2020).
\bibitem{page1954} Page, E. S. \emph{Biometrika} \textbf{41}, 100 (1954).
\bibitem{khan_delaylines} Khan, H. F., Kolhe, O. T., Habibi Matin, M. \&
Jayant, K. Neural delay lines synchronize inter-areal motor sequences.
\emph{Under review.}
\bibitem{gonzales2025} Gonzales, D. L., et al. \emph{Sci. Adv.} \textbf{11},
eadr4038 (2025).
\bibitem{jayant2019} Jayant, K., et al. \emph{Cell Rep.} \textbf{26}, 266
(2019).
\end{thebibliography}
\endgroup

\end{document}
